\documentclass[a4paper, 10pt, twoside]{article}

\usepackage[top=1in, bottom=1in, left=1in, right=1in]{geometry}
\usepackage[utf8]{inputenc}
\usepackage[spanish, es-ucroman, es-noquoting]{babel}
\usepackage{setspace}
\usepackage{fancyhdr}
\usepackage{lastpage}
\usepackage{amsmath}
\usepackage{amsfonts}
\usepackage{amsthm}
\usepackage{verbatim}
\usepackage{graphicx}
\usepackage{float}
\usepackage{enumitem} % Provee macro \setlist
\usepackage{tabularx}
\usepackage{multirow}
\usepackage{hyperref}
\usepackage{multicol}
\usepackage[toc, page]{appendix}


%%%%%%%%%% Configuración de Fancyhdr - Inicio %%%%%%%%%%
\pagestyle{fancy}
\thispagestyle{fancy}
\lhead{Trabajo Práctico 1 · Ingeniería de Software I}
\rhead{Delgado · Lovisolo · Petaccio · Requeni · Vita}
\renewcommand{\footrulewidth}{0.4pt}
\cfoot{\thepage /\pageref{LastPage}}

\fancypagestyle{caratula} {
   \fancyhf{}
   \cfoot{\thepage /\pageref{LastPage}}
   \renewcommand{\headrulewidth}{0pt}
   \renewcommand{\footrulewidth}{0pt}
}
%%%%%%%%%% Configuración de Fancyhdr - Fin %%%%%%%%%%


%%%%%%%%%% Miscelánea - Inicio %%%%%%%%%%
% Evita que el documento se estire verticalmente para ocupar el espacio vacío
% en cada página.
\raggedbottom

% Deshabilita sangría en la primer línea de un párrafo.
\setlength{\parindent}{0em}

% Separación entre párrafos.
\setlength{\parskip}{0.5em}

% Separación entre elementos de listas.
\setlist{itemsep=0.5em}

% Asigna la traducción de la palabra 'Appendices'.
\renewcommand{\appendixtocname}{Apéndices}
\renewcommand{\appendixpagename}{Apéndices}
%%%%%%%%%% Miscelánea - Fin %%%%%%%%%%


%%%%%%%%%% Insertar diagrama - Inicio %%%%%%%%%%
\newcommand{\diagrama}[1]{
  \includegraphics[width=16cm]{#1}
}
%%%%%%%%%% Insertar diagrama - Fin %%%%%%%%%%


\begin{document}


%%%%%%%%%%%%%%%%%%%%%%%%%%%%%%%%%%%%%%%%%%%%%%%%%%%%%%%%%%%%%%%%%%%%%%%%%%%%%%%
%% Carátula                                                                  %%
%%%%%%%%%%%%%%%%%%%%%%%%%%%%%%%%%%%%%%%%%%%%%%%%%%%%%%%%%%%%%%%%%%%%%%%%%%%%%%%


\thispagestyle{caratula}

\begin{center}

\includegraphics[height=2cm]{DC.png} 
\hfill
\includegraphics[height=2cm]{UBA.jpg} 

\vspace{2cm}

Departamento de Computación,\\
Facultad de Ciencias Exactas y Naturales,\\
Universidad de Buenos Aires

\vspace{4cm}

\begin{Huge}
Trabajo Práctico 1
\end{Huge}

\vspace{0.5cm}

\begin{Large}
Ingeniería de Software I
\end{Large}

\vspace{1cm}

Primer Cuatrimestre de 2014

\vspace{4cm}

\begin{tabular}{|c|c|c|}
\hline
Apellido y Nombre & LU & E-mail\\
\hline
Delgado Alejandro N.  & 601/11 & nahueldelgado@gmail.com\\
Lovisolo Leandro      & 645/11 & leandro@leandro.me\\
Petaccio Lautaro José & 443/11 & lausuper@gmail.com\\
Requeni Gastón        & 400/11 & grequeni@hotmail.com\\
Vita Sebastián        & 149/11 & sebastian\_vita@yahoo.com.ar\\
\hline
\end{tabular}

\end{center}

\newpage


%%%%%%%%%%%%%%%%%%%%%%%%%%%%%%%%%%%%%%%%%%%%%%%%%%%%%%%%%%%%%%%%%%%%%%%%%%%%%%%
%% Índice                                                                    %%
%%%%%%%%%%%%%%%%%%%%%%%%%%%%%%%%%%%%%%%%%%%%%%%%%%%%%%%%%%%%%%%%%%%%%%%%%%%%%%%


\tableofcontents

\newpage


%%%%%%%%%%%%%%%%%%%%%%%%%%%%%%%%%%%%%%%%%%%%%%%%%%%%%%%%%%%%%%%%%%%%%%%%%%%%%%%
%% Introducción                                                              %%
%%%%%%%%%%%%%%%%%%%%%%%%%%%%%%%%%%%%%%%%%%%%%%%%%%%%%%%%%%%%%%%%%%%%%%%%%%%%%%%


\section{Introducción}

En este trabajo práctico aplicamos varias técnicas de ingeniería de requerimientos para modelar un hipotético sistema de administración de la red de ciclovías de una ciudad. Además, proporcionamos una lista de escenarios informales que ejemplifican situaciones representativas del funcionamiento esperado.


%%%%%%%%%%%%%%%%%%%%%%%%%%%%%%%%%%%%%%%%%%%%%%%%%%%%%%%%%%%%%%%%%%%%%%%%%%%%%%%
%% Modelo de Jackson                                                         %%
%%%%%%%%%%%%%%%%%%%%%%%%%%%%%%%%%%%%%%%%%%%%%%%%%%%%%%%%%%%%%%%%%%%%%%%%%%%%%%%


\section{Modelo de Jackson}


\subsection{Mundo}

\begin{multicols}{2}
  \begin{itemize}
    \item Usuario retira bicicleta	
    \item Usuario usa la bicicleta más de una hora
    \item Usuario no consigue bicicleta
    \item Usuario sufre el robo de una bicicleta
    \item Usuario reporta al personal de la estación que su bicicleta fue robada	
    \item Usuario tiene accidente en bicicleta
    \item Usuarios salen de trabajar en una franja horaria reducida	
    \item Usuario espera bicicleta
    \item Usuario espera bicicleta por más de 40 minutos 
    \item Usuario roba identidad a otro
    \item Usuario penalizado retira bicicleta	
    \item Usuario sin registrar retira bicicleta
    \item Usuario retorna bicicleta tarde (más de 1hr) y se lo penaliza
    \item Usuario va a buscar una bici a estación del centro
    \item Usuario va a buscar una bici a estación periférica
    \item Usuario retorna bicicleta en horario 
    \item Personal de empresa de transporte trasladan bicicletas
    \item Personal de empresa de transporte tardan en trasladar bicicletas
    \item Personal de empresa de transporte lleva bicicletas a una estación
    \item Personal de empresa de transporte retira bicicletas de una estación
    \item Empresa de transporte recibe solicitud de transporte de bicis de estación A a B 
    \item Sistema pierde conexión
    \item Estacion pierde conexión
    \item Estación se queda sin bicicletas
    \item Personal de estación entrega una bicicleta 
    \item Personal de estación recibe una bicicleta
    \item Personal de estación rechaza una petición de retirar bicicleta debido suspensión por infracciones
    \item Personal de estación rechaza una petición de retirar bicicleta ya que el usuario ya tiene una bicicleta
    \item Personal de estación rechaza una petición de retirar bicicleta debido a que el usuario no es válido
    \item Personal de estación rechaza una petición de retirar una bicicleta ya que no hay stock
    \item Usuario se identifica con el personal de una estación
    \item Se cae la conexión de la estación
    \item Personal del gobierno de Mar Chiquita entrega nuevas bicicletas al personal de la empresa de transporte
  \end{itemize}
\end{multicols}


\subsection{Interfaz}

\begin{multicols}{2}
  \begin{itemize}
    \item Usuario se registra en el sistema
    \item Usuario consulta disponibilidad de bicicletas
    \item Usuario consulta su estado de penalización
    \item Personal de estación informa al sistema que una bicicleta fue robada
    \item Personal de estación informa al sistema el retiro de una bicicleta por parte de un usuario determinado
    \item Personal de estación informa al sistema la devolución de una bicicleta por parte de un usuario determinado
    \item Personal de estación solicita reposición de bicicletas
    \item Personal de estación autentifica un usuario en el sistema
    \item Personal de estación consulta con el sistema el retiro de una bicicleta por parte de un usuario determinado
    \item Sistema informa al personal de la estación si la petición de retirar bicicleta de un usuario es aceptada/rechazada
    \item Sistema informa al personal de la estación si la petición de autenticación de un usuario resulta correcta/incorrecta
    \item Sistema informa al usuario el resultado del registro
    \item Sistema informa al usuario el estado de penalización de este
    \item Sistema informa al usuario la disponibilidad de bicicletas
    \item Sistema envía ordenes de reposición de bicicletas al personal de la empresa de transporte
    \item Personal del gobierno de Mar Chiquita se autentica en el sistema
    \item Personal del gobierno de Mar Chiquita informa al sistema la disponibilidad de nuevas bicicletas
    \item Personal del gobierno de Mar Chiquita informa al sistema que existe una nueva estación

  \end{itemize}
\end{multicols}


\subsection{Sistema}

\begin{multicols}{2}
  \begin{itemize}
    \item Sistema checkea que un usuario registrado no vuelva a registrarse
    \item Sistema checkea si existen penalizaciones
    \item Sistema checkea si ya se sacó una bici
    \item Sistema penaliza usuarios por pérdia/robo de bicicleta
    \item Sistema calcula la penalización de los usuarios infractores basado en el tiempo excedido del límite de 1h
    \item Sistema registra usuario
    \item Sistema registra una nueva estacion
    \item Sistema registra nuevas bicicletas agregadas por el personal del estado
    \item Sistema penaliza usuario por pérdida de bicicleta
    \item Sistema utiliza un algoritmo de encriptación seguro para las contraseñas de los usuarios
    \item Sistema lleva registro de cantidad de bicis en cada estación 
    \item Sistema lleva registro de todos los usuarios existentes y si tienen o no una bici en su poder
    \item Sistema diferencia entre estaciones del centro y periféricas 
    \item Sistema trata de que las estaciones del centro tengan más bicis que las periféricas a todo momento (si se desbalancea mucho, repone automáticamente)
    \item Sistema le indica al personal de empresa de transporte cuantas bicicletas deben quitar de una estación y a donde entregarlas.
    \item Sistema calcula mediante un algoritmo que estación es la mas indicada para ceder bicicletas según la distancia entre las estaciones (?)
  \end{itemize}
\end{multicols}


%%%%%%%%%%%%%%%%%%%%%%%%%%%%%%%%%%%%%%%%%%%%%%%%%%%%%%%%%%%%%%%%%%%%%%%%%%%%%%%
%% Diagrama de Contexto                                                      %%
%%%%%%%%%%%%%%%%%%%%%%%%%%%%%%%%%%%%%%%%%%%%%%%%%%%%%%%%%%%%%%%%%%%%%%%%%%%%%%%


\section{Diagrama de Contexto}

\diagrama{diagrama-de-contexto.png}


%%%%%%%%%%%%%%%%%%%%%%%%%%%%%%%%%%%%%%%%%%%%%%%%%%%%%%%%%%%%%%%%%%%%%%%%%%%%%%%
%% Modelo de Objetivos                                                       %%
%%%%%%%%%%%%%%%%%%%%%%%%%%%%%%%%%%%%%%%%%%%%%%%%%%%%%%%%%%%%%%%%%%%%%%%%%%%%%%%


\section{Modelo de Objetivos}

\diagrama{modelo-de-objetivos.jpg}


%%%%%%%%%%%%%%%%%%%%%%%%%%%%%%%%%%%%%%%%%%%%%%%%%%%%%%%%%%%%%%%%%%%%%%%%%%%%%%%
%% Escenarios hipotéticos                                                    %%
%%%%%%%%%%%%%%%%%%%%%%%%%%%%%%%%%%%%%%%%%%%%%%%%%%%%%%%%%%%%%%%%%%%%%%%%%%%%%%%


\section{Escenarios Hipotéticos}


Los siguientes escenarios ilustran de manera informal situaciones representativas del funcionamiento esperado del sistema.


\subsection{Escenario 1}

Lorem ipsum dolor sit amet, consectetur adipiscing elit. Phasellus facilisis neque eget suscipit faucibus. Curabitur nunc ligula, molestie quis velit molestie, convallis molestie tellus. Suspendisse malesuada, justo euismod varius commodo, dolor lectus tristique justo, at malesuada velit eros non tellus. Sed convallis dolor sit amet nisi iaculis, a pharetra justo tempus. Cras sit amet urna in metus mattis blandit. Quisque laoreet vitae dolor ut pharetra. Proin ut vulputate leo, vel suscipit lacus. Nunc ultrices ante ut sem faucibus, ac viverra enim faucibus. Morbi feugiat tempus enim, dapibus lobortis tortor fringilla a. Praesent sem nibh, dignissim ac dictum vitae, convallis mollis magna. Ut tincidunt ipsum mi, at venenatis quam viverra nec. Integer justo nisl, auctor nec vestibulum non, interdum a mauris.

Etiam lacinia diam a ante sagittis cursus. Quisque in vestibulum purus. Duis eget vestibulum velit. Curabitur ornare mi a dapibus mattis. Vestibulum ac fermentum ipsum. Mauris pellentesque mattis risus, sit amet luctus mauris malesuada non. Sed vehicula pharetra neque, sagittis rhoncus est tempor in. Sed auctor consequat lectus ut congue. Vestibulum eget pharetra justo, consectetur porttitor tellus. Proin vel metus eu dolor ullamcorper ullamcorper. Sed nec nunc erat.


\subsection{Escenario 2}

Mauris vehicula sapien mauris, non lobortis lacus convallis sed. Donec sollicitudin ultrices dui. Nulla fermentum posuere odio, eget fermentum magna lacinia sit amet. Donec molestie tellus ac mi molestie, ac commodo massa laoreet. Nunc purus ante, consectetur et diam a, dapibus rutrum mi. Ut porttitor iaculis magna eget feugiat. Ut sit amet nulla aliquet, placerat odio vel, luctus libero. Etiam commodo egestas diam mollis laoreet. Fusce rutrum lectus ac augue sodales, eu tincidunt ligula varius. Pellentesque eget volutpat est. Aliquam sem purus, porta a dui a, rhoncus facilisis velit. Quisque at ullamcorper lectus. Ut in justo consectetur, scelerisque massa eu, faucibus ligula. Donec fermentum porta tellus, a venenatis quam varius quis. Nullam ac massa mattis, posuere neque at, egestas tellus.

Proin sit amet nulla tempor, cursus mi sed, tincidunt turpis. Nulla fermentum, arcu vitae imperdiet condimentum, nibh nisl euismod nibh, et suscipit turpis neque ut metus. Fusce vitae magna sem. Nullam aliquet, ipsum eu commodo gravida, leo odio porta turpis, sit amet pellentesque tellus felis non sapien. Phasellus pretium, felis a tempor convallis, justo turpis feugiat nibh, ac ornare lorem leo quis turpis. Nulla ipsum felis, scelerisque in arcu vitae, lobortis venenatis velit. Nam quis iaculis mi.


\subsection{Escenario 3}

Ut ullamcorper dignissim enim. Curabitur rhoncus orci ante, vitae iaculis lorem aliquam eu. Donec lobortis, purus sit amet ultrices molestie, felis dolor ultrices turpis, eu eleifend arcu turpis sed sem. Aliquam fringilla iaculis feugiat. Fusce vitae rhoncus felis. Aenean pellentesque dignissim lectus at aliquam. Duis condimentum aliquam auctor. Praesent porta scelerisque massa vitae sagittis. Praesent erat urna, ullamcorper ac nisl id, volutpat facilisis lorem. Nunc placerat dui posuere, faucibus ante eget, volutpat neque.

Vivamus neque nisl, congue non ornare non, pretium adipiscing nibh. Pellentesque rhoncus tincidunt commodo. Morbi nulla libero, egestas in congue eu, tempor nec leo. Donec at ornare ante. Duis vulputate pellentesque risus sed imperdiet. Aenean tempus id risus vitae consectetur. Curabitur consequat varius orci, id accumsan quam laoreet ut. Sed dui velit, semper ac posuere eu, consequat eget dolor. Integer ut diam pellentesque, tristique nisl eget, posuere elit. Etiam vestibulum arcu at lacus lobortis molestie. Suspendisse bibendum ante non dolor mollis, at euismod enim blandit. Nullam sagittis velit dignissim, pretium libero volutpat, adipiscing tortor


\end{document}